%
% repl.tex
%
% Z Specification for Biff REPL Lifecycle
% Formal model of session lifecycle, command dispatch, prompt display
% coordination, and notification state for the interactive REPL.
%

\documentclass[a4paper,10pt,fleqn]{article}
\usepackage[margin=1in]{geometry}
\usepackage{fuzz}
\usepackage[colorlinks=true,linkcolor=blue,citecolor=blue,urlcolor=blue]{hyperref}

\begin{document}

\title{Biff REPL Lifecycle: A Z Specification}
\author{Formal Model of Session, Dispatch, Prompt, and Notification State}
\date{March 2026}
\hypersetup{
  pdftitle={Biff REPL Lifecycle: A Z Specification},
  pdfauthor={punt-labs},
  pdfsubject={Z Specification},
  pdfcreator={fuzz/probcli}
}
\maketitle

\tableofcontents
\newpage

%%----------------------------------------------------------------------
\section{Introduction}
%%----------------------------------------------------------------------

The biff REPL (\texttt{biff} with no arguments) provides an interactive
command loop with a proper session lifecycle that mirrors the MCP
server.  This specification models four concurrent state domains and
their interactions:

\begin{enumerate}
\item \textbf{Session lifecycle} --- connect to NATS, register in KV
  with an auto-assigned \texttt{ttyN} name, write wtmp login, run
  heartbeat, and clean up on exit (wtmp logout, KV delete, disconnect).

\item \textbf{Command dispatch} --- the stdin reader thread puts lines
  into an asyncio queue; the event loop drains the queue and dispatches
  commands via \texttt{shlex}-based parsing.  Commands return a
  \texttt{CommandResult} with text and optional JSON.  The special
  commands \texttt{exit}/\texttt{quit} terminate the loop; \texttt{talk}
  enters a modal sub-loop (modelled in \texttt{docs/talk.tex}).

\item \textbf{Prompt display} --- a threading gate (prompt gate)
  synchronises the stdin reader thread and the async event loop.
  The thread waits for the gate before calling \texttt{input(prompt)};
  the async loop sets the gate after finishing command output.  This
  prevents the prompt from interleaving with output.

\item \textbf{Notification state} --- \texttt{NotifyState} tracks the
  last-known unread count and wall fingerprint.  On each poll (2-second
  timeout or NATS notification wake), the current values are compared
  with the snapshot.  Changes produce inline banner lines.  After the
  user's own command, \texttt{sync()} updates the snapshot without
  emitting banners (self-notification prevention).
\end{enumerate}

The prompt gate is the critical coordination mechanism.  It is the
formal equivalent of a semaphore controlling terminal write access
between the stdin reader thread and the async event loop.

%%----------------------------------------------------------------------
\section{Basic Types}
%%----------------------------------------------------------------------

\begin{zed}
  [USER, TTY, SESSIONKEY, CMDLINE, CMDOUTPUT, WALLKEY]
\end{zed}

$USER$ identifies a biff user.  $TTY$ is a session identifier.
$SESSIONKEY$ is a composite \texttt{user:tty} key.  $CMDLINE$ is
a raw input line from the user.  $CMDOUTPUT$ is the text produced
by a command.  $WALLKEY$ is a wall fingerprint (text + posted\_at).

%%----------------------------------------------------------------------
\section{Free Types}
%%----------------------------------------------------------------------

\begin{zed}
  ZBOOL ::= ztrue | zfalse
\end{zed}

\begin{zed}
  SessionPhase ::= spDisconnected | spConnected | spShuttingDown
\end{zed}

The session connection lifecycle: initially disconnected, connected
after \texttt{cli\_session} enters, and shutting down during the
\texttt{finally} block.

\begin{zed}
  ReplMode ::= rmNormal | rmTalk
\end{zed}

The REPL is either in normal command mode or in the modal talk
sub-loop.  The talk sub-loop is modelled in \texttt{docs/talk.tex};
this specification models the mode transition.

\begin{zed}
  DispatchResult ::= drOutput | drExit | drTalk | drError | drEmpty
\end{zed}

The five outcomes of dispatching a command line:
\begin{itemize}
\item $drOutput$ --- command produced output to display.
\item $drExit$ --- user typed \texttt{exit} or \texttt{quit}.
\item $drTalk$ --- user typed \texttt{talk @user}; enter talk mode.
\item $drError$ --- command raised an error.
\item $drEmpty$ --- blank line; no action.
\end{itemize}

%%----------------------------------------------------------------------
\section{Global Constants}
%%----------------------------------------------------------------------

\begin{axdef}
  maxUnread : \nat \\
  heartbeatInterval : \nat \\
  pollTimeout : \nat
\where
  maxUnread = 100 \\
  heartbeatInterval = 60 \\
  pollTimeout = 2
\end{axdef}

$maxUnread$ caps the unread count (DES-015).  $heartbeatInterval$
is the KV refresh period in seconds.  $pollTimeout$ is the
notification poll cycle in seconds.

%%----------------------------------------------------------------------
\section{REPL State}
%%----------------------------------------------------------------------

All state is flattened into a single schema for ProB compatibility.

\begin{schema}{ReplState}
  sessionPhase : SessionPhase \\
  replMode : ReplMode \\
  myUser : USER \\
  myTty : TTY \\
  myKey : SESSIONKEY \\
  registered : ZBOOL \\
  promptGateOpen : ZBOOL \\
  lastUnread : \nat \\
  lastWallKey : WALLKEY \\
  currentUnread : \nat \\
  currentWallKey : WALLKEY \\
  pendingOutput : ZBOOL \\
  inputAvailable : ZBOOL
\where
  lastUnread \leq maxUnread \\
  currentUnread \leq maxUnread \\
  sessionPhase = spDisconnected \implies registered = zfalse \\
  sessionPhase = spDisconnected \implies promptGateOpen = zfalse \\
  replMode = rmTalk \implies sessionPhase = spConnected \\
  pendingOutput = ztrue \implies promptGateOpen = zfalse
\end{schema}

The key invariants:
\begin{itemize}
\item A disconnected session is never registered and has no open
  prompt gate.
\item Talk mode requires an active connection.
\item When output is pending (the async loop is printing), the prompt
  gate is closed --- the stdin thread cannot print a prompt.
\end{itemize}

%%----------------------------------------------------------------------
\section{Initialization}
%%----------------------------------------------------------------------

\begin{schema}{Init}
  ReplState'
\where
  sessionPhase' = spDisconnected \\
  replMode' = rmNormal \\
  registered' = zfalse \\
  promptGateOpen' = zfalse \\
  lastUnread' = 0 \\
  currentUnread' = 0 \\
  pendingOutput' = zfalse \\
  inputAvailable' = zfalse
\end{schema}

The initial state before \texttt{cli\_session} connects.  Wall keys
and identity ($myUser$, $myTty$, $myKey$) are unconstrained
(assigned during connect).

%%----------------------------------------------------------------------
\section{Session Lifecycle Operations}
%%----------------------------------------------------------------------

%%----------------------------------------------------------------------
\subsection{Connect}
%%----------------------------------------------------------------------

\texttt{cli\_session} enters: connect to NATS, register session in
KV, auto-assign ttyN, write wtmp login, start heartbeat, seed
notification state.

\begin{schema}{Connect}
  \Delta ReplState
\where
  sessionPhase = spDisconnected \\
  sessionPhase' = spConnected \\
  registered' = ztrue \\
  promptGateOpen' = ztrue \\
  replMode' = rmNormal \\
  pendingOutput' = zfalse \\
  inputAvailable' = zfalse \\
  lastUnread' = currentUnread' \\
  myUser' = myUser \\
  myTty' = myTty \\
  myKey' = myKey
\end{schema}

After connect, the prompt gate opens to allow the first prompt.
Notification state is seeded ($lastUnread' = currentUnread'$) so
pre-existing unread messages do not trigger a banner.

%%----------------------------------------------------------------------
\subsection{Heartbeat}
%%----------------------------------------------------------------------

Periodic KV refresh.  No state change visible to the REPL; modelled
as a query to show it preserves all invariants.

\begin{schema}{Heartbeat}
  \Xi ReplState
\where
  sessionPhase = spConnected
\end{schema}

%%----------------------------------------------------------------------
\subsection{Disconnect}
%%----------------------------------------------------------------------

\texttt{cli\_session} exits: write wtmp logout, delete KV entry,
close NATS connection.

\begin{schema}{Disconnect}
  \Delta ReplState
\where
  sessionPhase = spConnected \\
  sessionPhase' = spShuttingDown \\
  registered' = zfalse \\
  promptGateOpen' = zfalse \\
  replMode' = rmNormal \\
  pendingOutput' = zfalse \\
  inputAvailable' = zfalse \\
  lastUnread' = lastUnread \\
  lastWallKey' = lastWallKey \\
  currentUnread' = currentUnread \\
  currentWallKey' = currentWallKey \\
  myUser' = myUser \\
  myTty' = myTty \\
  myKey' = myKey
\end{schema}

%%----------------------------------------------------------------------
\section{Input/Dispatch/Output Cycle}
%%----------------------------------------------------------------------

%%----------------------------------------------------------------------
\subsection{InputReady}
%%----------------------------------------------------------------------

The stdin reader thread has put a line into the asyncio queue.
The prompt gate was open (thread was allowed to read); now it
closes because the async loop needs exclusive terminal access to
print output.

\begin{schema}{InputReady}
  \Delta ReplState
\where
  sessionPhase = spConnected \\
  replMode = rmNormal \\
  promptGateOpen = ztrue \\
  inputAvailable = zfalse \\
  inputAvailable' = ztrue \\
  promptGateOpen' = zfalse \\
  sessionPhase' = sessionPhase \\
  replMode' = replMode \\
  registered' = registered \\
  pendingOutput' = pendingOutput \\
  lastUnread' = lastUnread \\
  lastWallKey' = lastWallKey \\
  currentUnread' = currentUnread \\
  currentWallKey' = currentWallKey \\
  myUser' = myUser \\
  myTty' = myTty \\
  myKey' = myKey
\end{schema}

%%----------------------------------------------------------------------
\subsection{DispatchCommand}
%%----------------------------------------------------------------------

The async loop takes a line from the queue and dispatches it.
Output is generated; the input is consumed.

\begin{schema}{DispatchCommand}
  \Delta ReplState \\
  result? : DispatchResult
\where
  sessionPhase = spConnected \\
  replMode = rmNormal \\
  inputAvailable = ztrue \\
  result? = drOutput \\
  inputAvailable' = zfalse \\
  pendingOutput' = ztrue \\
  sessionPhase' = sessionPhase \\
  replMode' = replMode \\
  registered' = registered \\
  promptGateOpen' = promptGateOpen \\
  lastUnread' = lastUnread \\
  lastWallKey' = lastWallKey \\
  currentUnread' = currentUnread \\
  currentWallKey' = currentWallKey \\
  myUser' = myUser \\
  myTty' = myTty \\
  myKey' = myKey
\end{schema}

%%----------------------------------------------------------------------
\subsection{DispatchExit}
%%----------------------------------------------------------------------

User typed \texttt{exit} or \texttt{quit}.  Triggers disconnect.

\begin{schema}{DispatchExit}
  \Delta ReplState
\where
  sessionPhase = spConnected \\
  replMode = rmNormal \\
  inputAvailable = ztrue \\
  inputAvailable' = zfalse \\
  sessionPhase' = spShuttingDown \\
  promptGateOpen' = zfalse \\
  replMode' = rmNormal \\
  registered' = registered \\
  pendingOutput' = zfalse \\
  lastUnread' = lastUnread \\
  lastWallKey' = lastWallKey \\
  currentUnread' = currentUnread \\
  currentWallKey' = currentWallKey \\
  myUser' = myUser \\
  myTty' = myTty \\
  myKey' = myKey
\end{schema}

%%----------------------------------------------------------------------
\subsection{DispatchEmpty}
%%----------------------------------------------------------------------

Blank line.  No output, re-open the prompt gate immediately.

\begin{schema}{DispatchEmpty}
  \Delta ReplState
\where
  sessionPhase = spConnected \\
  replMode = rmNormal \\
  inputAvailable = ztrue \\
  inputAvailable' = zfalse \\
  promptGateOpen' = ztrue \\
  pendingOutput' = zfalse \\
  sessionPhase' = sessionPhase \\
  replMode' = replMode \\
  registered' = registered \\
  lastUnread' = lastUnread \\
  lastWallKey' = lastWallKey \\
  currentUnread' = currentUnread \\
  currentWallKey' = currentWallKey \\
  myUser' = myUser \\
  myTty' = myTty \\
  myKey' = myKey
\end{schema}

%%----------------------------------------------------------------------
\subsection{DispatchError}
%%----------------------------------------------------------------------

Command raised a \texttt{ValueError}.  Error printed to stderr,
prompt gate re-opened.

\begin{schema}{DispatchError}
  \Delta ReplState
\where
  sessionPhase = spConnected \\
  replMode = rmNormal \\
  inputAvailable = ztrue \\
  inputAvailable' = zfalse \\
  promptGateOpen' = ztrue \\
  pendingOutput' = zfalse \\
  sessionPhase' = sessionPhase \\
  replMode' = replMode \\
  registered' = registered \\
  lastUnread' = lastUnread \\
  lastWallKey' = lastWallKey \\
  currentUnread' = currentUnread \\
  currentWallKey' = currentWallKey \\
  myUser' = myUser \\
  myTty' = myTty \\
  myKey' = myKey
\end{schema}

%%----------------------------------------------------------------------
\subsection{OutputComplete}
%%----------------------------------------------------------------------

Command output has been printed and notification state synced.
Re-open the prompt gate so the stdin thread can print the next
prompt.

\begin{schema}{OutputComplete}
  \Delta ReplState
\where
  sessionPhase = spConnected \\
  pendingOutput = ztrue \\
  pendingOutput' = zfalse \\
  promptGateOpen' = ztrue \\
  lastUnread' = currentUnread \\
  lastWallKey' = currentWallKey \\
  sessionPhase' = sessionPhase \\
  replMode' = replMode \\
  registered' = registered \\
  inputAvailable' = inputAvailable \\
  currentUnread' = currentUnread \\
  currentWallKey' = currentWallKey \\
  myUser' = myUser \\
  myTty' = myTty \\
  myKey' = myKey
\end{schema}

The \texttt{sync()} call happens here: $lastUnread' = currentUnread$
and $lastWallKey' = currentWallKey$.  This prevents the next poll
from generating a self-notification for changes the user's own
command just made.

%%----------------------------------------------------------------------
\subsection{EnterTalkMode}
%%----------------------------------------------------------------------

User typed \texttt{talk @user}.  REPL transitions to talk mode
(modelled in \texttt{docs/talk.tex}).

\begin{schema}{EnterTalkMode}
  \Delta ReplState
\where
  sessionPhase = spConnected \\
  replMode = rmNormal \\
  inputAvailable = ztrue \\
  replMode' = rmTalk \\
  inputAvailable' = zfalse \\
  promptGateOpen' = zfalse \\
  pendingOutput' = zfalse \\
  sessionPhase' = sessionPhase \\
  registered' = registered \\
  lastUnread' = lastUnread \\
  lastWallKey' = lastWallKey \\
  currentUnread' = currentUnread \\
  currentWallKey' = currentWallKey \\
  myUser' = myUser \\
  myTty' = myTty \\
  myKey' = myKey
\end{schema}

%%----------------------------------------------------------------------
\subsection{ExitTalkMode}
%%----------------------------------------------------------------------

Talk sub-loop ended (local or remote hangup).  Return to normal
mode, re-open prompt gate.

\begin{schema}{ExitTalkMode}
  \Delta ReplState
\where
  sessionPhase = spConnected \\
  replMode = rmTalk \\
  replMode' = rmNormal \\
  promptGateOpen' = ztrue \\
  pendingOutput' = zfalse \\
  inputAvailable' = zfalse \\
  sessionPhase' = sessionPhase \\
  registered' = registered \\
  lastUnread' = lastUnread \\
  lastWallKey' = lastWallKey \\
  currentUnread' = currentUnread \\
  currentWallKey' = currentWallKey \\
  myUser' = myUser \\
  myTty' = myTty \\
  myKey' = myKey
\end{schema}

%%----------------------------------------------------------------------
\section{Notification Operations}
%%----------------------------------------------------------------------

%%----------------------------------------------------------------------
\subsection{PollNoChange}
%%----------------------------------------------------------------------

2-second timeout fires.  Current unread count and wall key are the
same as the snapshot.  No banners emitted.

\begin{schema}{PollNoChange}
  \Xi ReplState
\where
  sessionPhase = spConnected \\
  replMode = rmNormal \\
  currentUnread = lastUnread \\
  currentWallKey = lastWallKey
\end{schema}

%%----------------------------------------------------------------------
\subsection{PollNewMessages}
%%----------------------------------------------------------------------

Unread count increased since last snapshot.  A banner is emitted
inline (clear line, print, reshow prompt).  Snapshot updated.

\begin{schema}{PollNewMessages}
  \Delta ReplState
\where
  sessionPhase = spConnected \\
  replMode = rmNormal \\
  currentUnread > lastUnread \\
  lastUnread' = currentUnread \\
  lastWallKey' = lastWallKey \\
  sessionPhase' = sessionPhase \\
  replMode' = replMode \\
  registered' = registered \\
  promptGateOpen' = promptGateOpen \\
  pendingOutput' = pendingOutput \\
  inputAvailable' = inputAvailable \\
  currentUnread' = currentUnread \\
  currentWallKey' = currentWallKey \\
  myUser' = myUser \\
  myTty' = myTty \\
  myKey' = myKey
\end{schema}

%%----------------------------------------------------------------------
\subsection{PollWallChanged}
%%----------------------------------------------------------------------

Wall fingerprint changed (new wall posted, or wall expired/cleared).
A banner is emitted inline.  Snapshot updated.

\begin{schema}{PollWallChanged}
  \Delta ReplState
\where
  sessionPhase = spConnected \\
  replMode = rmNormal \\
  currentWallKey \neq lastWallKey \\
  lastWallKey' = currentWallKey \\
  lastUnread' = lastUnread \\
  sessionPhase' = sessionPhase \\
  replMode' = replMode \\
  registered' = registered \\
  promptGateOpen' = promptGateOpen \\
  pendingOutput' = pendingOutput \\
  inputAvailable' = inputAvailable \\
  currentUnread' = currentUnread \\
  currentWallKey' = currentWallKey \\
  myUser' = myUser \\
  myTty' = myTty \\
  myKey' = myKey
\end{schema}

%%----------------------------------------------------------------------
\subsection{ExternalStateChange}
%%----------------------------------------------------------------------

An external event changes the unread count or wall key (message
arrival, wall post by another user).  This is the ``world'' updating
the current values that the next poll will compare against the
snapshot.

\begin{schema}{ExternalStateChange}
  \Delta ReplState \\
  newUnread? : \nat \\
  newWallKey? : WALLKEY
\where
  newUnread? \leq maxUnread \\
  currentUnread' = newUnread? \\
  currentWallKey' = newWallKey? \\
  sessionPhase' = sessionPhase \\
  replMode' = replMode \\
  registered' = registered \\
  promptGateOpen' = promptGateOpen \\
  pendingOutput' = pendingOutput \\
  inputAvailable' = inputAvailable \\
  lastUnread' = lastUnread \\
  lastWallKey' = lastWallKey \\
  myUser' = myUser \\
  myTty' = myTty \\
  myKey' = myKey
\end{schema}

%%----------------------------------------------------------------------
\section{System Invariants}
%%----------------------------------------------------------------------

The following properties must hold in all reachable states:

\begin{enumerate}
\item \textbf{Prompt gate closed during output.}  When
  $pendingOutput = ztrue$, $promptGateOpen = zfalse$.  The stdin
  thread never prints a prompt while the async loop is writing
  command output.

\item \textbf{Disconnected implies no registration.}  When
  $sessionPhase = spDisconnected$, $registered = zfalse$.

\item \textbf{Talk requires connection.}  When $replMode = rmTalk$,
  $sessionPhase = spConnected$.

\item \textbf{Notification snapshot bounded.}  $lastUnread$ and
  $currentUnread$ are both at most $maxUnread$.

\item \textbf{Self-notification prevention.}  After
  $OutputComplete$, $lastUnread = currentUnread$ and
  $lastWallKey = currentWallKey$.  The next $PollNoChange$ will
  emit no banners.

\item \textbf{Prompt gate lifecycle.}  The prompt gate is opened
  exactly once per command cycle: either in $OutputComplete$ (after
  successful command), $DispatchEmpty$ (blank line), $DispatchError$
  (error), or $ExitTalkMode$ (returning from talk).  It is closed
  by $InputReady$ when the stdin thread delivers a line.
\end{enumerate}

\end{document}
